% Идея метода отсечения: объемлющий многогранник, решение на нём,
% отсекающая плоскость, уменьшенный многогранник.
\begin{figure}[htbp]
\centering
\begin{tikzpicture}[scale=1.0, >={Stealth[length=2mm]}]

% объемлющий многогранник P_0
\draw[blue!55!black, thick]
  (-0.6,0.9) -- (1.3,-0.5) -- (4.0,-0.3) -- (5.0,1.6)
  -- (3.4,3.3) -- (0.6,3.0) -- cycle;
\node[blue!55!black, font=\footnotesize] at (4.5,2.9) {$P_0$};

% само допустимое множество S (гладкое, внутри)
\fill[blue!12] (2.2,1.4) ellipse (1.55 and 1.05);
\draw[blue!55!black, thick] (2.2,1.4) ellipse (1.55 and 1.05);
\node[blue!55!black, font=\footnotesize] at (2.2,1.4) {$S$};

% решение на P_0 — вершина, лежащая вне S
\coordinate (v) at (5.0,1.6);
\fill[orange!75!black] (v) circle (2pt);
\node[font=\scriptsize, anchor=west] at ($(v)+(0.12,0.05)$) {$x^{(0)} \notin S$};

% отсекающая плоскость: отрезает вершину, не задевая S
\draw[orange!75!black, very thick] (4.15,-0.22) -- (3.95,3.15);
\node[orange!60!black, font=\scriptsize, anchor=south west, rotate=87]
  at (4.05,0.1) {отсекающая плоскость};

% заштрихованная отсекаемая часть
\fill[orange!15] (4.15,-0.22) -- (4.0,-0.3) -- (5.0,1.6) -- (3.4,3.3)
                 -- (3.97,3.16) -- cycle;

\end{tikzpicture}
\caption{Идея метода отсечения. Допустимое множество $S$ заменяется объемлющим
многогранником $P_0$, на котором задача решается просто. Если найденный оптимум
$x^{(0)}$ не попал в $S$, через него проводится гиперплоскость, отсекающая
$x^{(0)}$, но не задевающая ни одной точки $S$. Получается
$P_1 \subset P_0$, по-прежнему содержащее $S$, и процедура повторяется.}
\label{fig:cutting-plane}
\end{figure}
